% Generated by scripts/build_reports.py; edit reviewed Markdown inputs.
% Compile with XeLaTeX.
% Options for packages loaded elsewhere
\PassOptionsToPackage{unicode}{hyperref}
\PassOptionsToPackage{hyphens}{url}
\PassOptionsToPackage{dvipsnames,svgnames,x11names}{xcolor}
\documentclass[
  11pt,
]{article}
\usepackage{xcolor}
\usepackage[margin=1in]{geometry}
\usepackage{amsmath,amssymb}
\setcounter{secnumdepth}{-\maxdimen} % remove section numbering
\usepackage{iftex}
\ifPDFTeX
  \usepackage[T1]{fontenc}
  \usepackage[utf8]{inputenc}
  \usepackage{textcomp} % provide euro and other symbols
\else % if luatex or xetex
  \usepackage{unicode-math} % this also loads fontspec
  \defaultfontfeatures{Scale=MatchLowercase}
  \defaultfontfeatures[\rmfamily]{Ligatures=TeX,Scale=1}
\fi
\usepackage{lmodern}
\ifPDFTeX\else
  % xetex/luatex font selection
\fi
% Use upquote if available, for straight quotes in verbatim environments
\IfFileExists{upquote.sty}{\usepackage{upquote}}{}
\IfFileExists{microtype.sty}{% use microtype if available
  \usepackage[]{microtype}
  \UseMicrotypeSet[protrusion]{basicmath} % disable protrusion for tt fonts
}{}
\makeatletter
\@ifundefined{KOMAClassName}{% if non-KOMA class
  \IfFileExists{parskip.sty}{%
    \usepackage{parskip}
  }{% else
    \setlength{\parindent}{0pt}
    \setlength{\parskip}{6pt plus 2pt minus 1pt}}
}{% if KOMA class
  \KOMAoptions{parskip=half}}
\makeatother
\usepackage{longtable,booktabs,array}
\usepackage{caption}
\captionsetup[table]{skip=6pt}
\newcounter{none} % for unnumbered tables
\usepackage{calc} % for calculating minipage widths
% Correct order of tables after \paragraph or \subparagraph
\usepackage{etoolbox}
\makeatletter
\patchcmd\longtable{\par}{\if@noskipsec\mbox{}\fi\par}{}{}
\makeatother
% Allow footnotes in longtable head/foot
\IfFileExists{footnotehyper.sty}{\usepackage{footnotehyper}}{\usepackage{footnote}}
\makesavenoteenv{longtable}
\setlength{\emergencystretch}{3em} % prevent overfull lines
\providecommand{\tightlist}{%
  \setlength{\itemsep}{0pt}\setlength{\parskip}{0pt}}
\usepackage{amsmath,amssymb}
\usepackage{xurl}
\usepackage{fancyhdr}
\pagestyle{fancy}
\fancyhf{}
\fancyhead[L]{\small Pure Thought | Revised research specifications}
\fancyfoot[L]{\small 2026-09-17}
\fancyfoot[R]{\small\thepage}
\setlength{\headheight}{15pt}
\setlength{\emergencystretch}{3em}
\allowdisplaybreaks
\usepackage{bookmark}
\IfFileExists{xurl.sty}{\usepackage{xurl}}{} % add URL line breaks if available
\urlstyle{same}
% fallback for those not using the hyperref driver hyperxmp:
\makeatletter
\@ifundefined{xmpquote}{\newcommand{\xmpquote}[1]{#1}}{}
\makeatother
\hypersetup{
  pdftitle={08. Exact modular-data and symmetry-anomaly consistency tests},
  pdfauthor={Pure Thought research portfolio},
  colorlinks=true,
  linkcolor={blue},
  filecolor={Maroon},
  citecolor={Blue},
  urlcolor={blue},
  pdfcreator={LaTeX via pandoc}}

\title{08. Exact modular-data and symmetry-anomaly consistency tests}
\author{Pure Thought research portfolio}
\date{Revised 17 September 2026}

\begin{document}
\maketitle
\thispagestyle{empty}
\clearpage

\section{08. Exact modular-data and symmetry-anomaly consistency tests}

\textbf{Revised specification --- 2026-09-11.}
\href{https://github.com/rick-stevens-ai/Pure-Thought/blob/main/reviews/08-Swampland-Modularity-Symmetries.md}{Detailed
critique} \ensuremath{\cdot}
\href{https://github.com/rick-stevens-ai/Pure-Thought/blob/main/archive/original-PRDs/08-Swampland-Modularity-Symmetries.md}{Original
description} \ensuremath{\cdot}
\href{https://github.com/rick-stevens-ai/Pure-Thought/blob/main/reviews/EVIDENCE-STANDARD.md}{Evidence
standard}

\subsection{Plain-language guide}

\subsubsection{The problem in everyday terms}

A proposed quantum theory comes with mathematical rules for combining
excitations and transforming its descriptions. These rules should fit
together like the pieces of a puzzle. This problem builds exact checks
for that fit and then asks a second question: can a chosen global
symmetry be made into a gauge symmetry? A theory can be consistent even
when this second operation is obstructed, so the two verdicts must
remain separate.

\subsubsection{Key terms}

\begin{itemize}
\item
  \textbf{Rational CFT} --- A conformal field theory with finitely many
  basic representation types relative to a specified chiral algebra.
\item
  \textbf{Character} --- A generating function counting states within
  one representation of that algebra.
\item
  \textbf{Modular S and T matrices} --- Matrices describing how the
  characters change under two basic transformations of the torus.
\item
  \textbf{Fusion rules} --- Rules specifying which excitation types can
  occur when two types are combined.
\item
  \textbf{Global symmetry} --- A transformation applied consistently
  throughout a system that leaves its laws unchanged.
\item
  \textbf{Gauging} --- Promoting an appropriate symmetry to a local
  redundancy and adding the corresponding gauge structure.
\item
  \textbf{Anomaly} --- An obstruction to implementing a classical
  symmetry or a proposed gauging consistently in the quantum setting; an
  't Hooft anomaly can exist in a valid theory.
\item
  \textbf{Cocycle} --- Algebraic data satisfying a compatibility
  identity, used here to encode an anomaly.
\item
  \textbf{Bulk and boundary} --- The interior gravitational theory and
  its lower-dimensional boundary description; symmetries can have
  different interpretations in the two.
\end{itemize}

\subsubsection{Why this matters}

Exact consistency checks can detect false candidate data before
researchers invest in a full construction. They also clarify which
symmetries can be gauged and which require additional structure, such as
anomaly cancellation or inflow from another system. This is useful both
for studying quantum field theories and for interpreting holographic
gravity proposals.

\subsubsection{What progress would mean}

The first contribution would be a dependable checker with known passing
and failing examples. A failure would identify a particular algebraic
inconsistency or gauging obstruction, rather than issuing an unsupported
verdict about all of quantum gravity.

\subsection{Core question}

For supplied rational-CFT modular data and a specified finite internal
symmetry, which algebraic consistency tests can be certified, and which
failures obstruct a proposed gauging?

\subsection{Scope and assumptions}

Use a bosonic rational CFT with a declared chiral algebra and finitely
many characters. Input exact algebraic S,T matrices, vacuum label, and a
nonnegative integer multiplicity matrix M. Start with the Ising model.
An optional finite 0-form group G requires separate
defect/twisted-sector data and an anomaly cocycle; these cannot be
inferred from the untwisted partition function alone.

The result concerns the supplied CFT data and gauging problem. A general
bulk quantum-gravity compatibility verdict is outside the baseline.

\subsection{Mathematical target}

Check S\ensuremath{\dagger}S=I, symmetry of S where assumed,
S\ensuremath{{}^{2}}=C, (ST)\ensuremath{{}^{3}}=C in the declared T
convention, M\ensuremath{{}_{00}}=1, MS=SM, MT=TM, and
M\ensuremath{{}_{ij}}\ensuremath{\in}Z\ensuremath{\ge}0. Check Verlinde
coefficients \ensuremath{\Sigma}\_a S\_ia S\_ja S*\_ka/S\_0a are
nonnegative integers.

For a finite symmetry cocycle
\ensuremath{\omega}\ensuremath{\in}Z\ensuremath{{}^{3}}(G,U(1)), check
the cocycle identity exactly. To certify anomaly-free gauging in the
stated setting, supply an explicit 2-cochain \ensuremath{\beta} with
\ensuremath{\delta}\ensuremath{\beta}=\ensuremath{\omega}, or prove its
nonexistence in a complete cohomology computation. Distinguish this
obstruction from inconsistency of an anomalous but valid boundary
theory.

\subsection{Required outputs and proof obligations}

Return pass/fail/unresolved for each named axiom, with exact arithmetic
traces. Include the number field and embeddings. Even passing every
modular-data test gives necessary algebraic consistency only; full CFT
sewing and realizability are separate obligations.

\subsection{Validation and rejection controls}

Verify Ising S entries including 1/\ensuremath{\surd}2 in algebraic
arithmetic. Reject altered modular relations, noninteger multiplicities,
and a failed cocycle identity. Accept the existence of a boundary global
symmetry rather than flagging it as automatically inconsistent.
Cross-check a small cyclic-group cohomology example.

\subsection{Milestones and research extension}

Baseline: exact modular-data checker on Ising and one second supplied
model. Strong: one finite-group twisted-sector/gauging analysis.
Research extension: classify a explicitly finite family of data or
derive a conditional bulk consequence using a stated holographic
theorem.

Milestones are acceptance gates, not calendar promises. Before a run,
freeze the concrete instance/domain, tolerance, compute budget, and
permitted dependencies. A bounded search that fails to find a result
must preserve an unresolved outcome.

\subsection{Starting instruction}

Work on the baseline above. First make every input and convention
explicit, then validate the reference case. Build the checker alongside
the calculation. Label numerical evidence, finite-model certificates,
and continuum theorems separately; return the strongest justified result
and identify the exact unmet proof obligation.

\subsection{Primary references and their role}

\begin{itemize}
\tightlist
\item
  \href{https://arxiv.org/abs/1810.05338}{Harlow and Ooguri, Symmetries
  in quantum field theory and quantum gravity} --- Bulk gauge versus
  boundary global symmetry, including higher-form extensions.
\end{itemize}

References support definitions or limitations; the scope choices and
proposed milestones are editorial recommendations, not claims that these
results have already been obtained.

\subsection{Displayed formulation and claim boundary}

\href{https://github.com/rick-stevens-ai/Pure-Thought/blob/main/PRDs/08-Swampland-Modularity-Symmetries.md}{Canonical
specification}.

\[S^\dagger S=I,\quad S^2=C,\quad(ST)^3=C,\quad
MS=SM,\quad MT=TM,\quad M_{00}=1,\quad M_{ij}\in\mathbb Z_{\ge0}.\]

\[N_{ij}^{\ k}=\sum_a\frac{S_{ia}S_{ja}S_{ka}^*}{S_{0a}}
\in\mathbb Z_{\ge0},\qquad \omega\in Z^3(G,U(1)),\quad\delta\beta=\omega.\]

Fix the \(T\) convention and any assumed symmetry of \(S\). An explicit
cochain \(\beta\) trivializes the specified anomaly; a complete
cohomology calculation may prove no such cochain exists. This gauging
obstruction does not make an anomalous boundary theory inconsistent, and
untwisted modular data alone do not supply defect data.

\subsection{Detailed review and rationale}

\href{https://github.com/rick-stevens-ai/Pure-Thought/blob/main/PRDs/08-Swampland-Modularity-Symmetries.md}{Revised
problem} \ensuremath{\cdot}
\href{https://github.com/rick-stevens-ai/Pure-Thought/blob/main/archive/original-PRDs/08-Swampland-Modularity-Symmetries.md}{Preserved
original} \ensuremath{\cdot}
\href{https://github.com/rick-stevens-ai/Pure-Thought/blob/f854d4e12b55c1f5b5c2a964b75f2d23496d981e/PRDs/08-Swampland-Modularity-Symmetries.md}{Pinned
upstream source}

Line references below refer to the preserved original at commit
\texttt{f854d4e12b55}, not the rewritten file. ``Blocking'' means the
original issue can invalidate the scientific target or claimed
certificate; ``major'' means it materially changes scope, inference, or
acceptance criteria. These are critiques of specifications, not results
of running the embedded research-code sketches.

\subsubsection{Detailed findings}

\begin{enumerate}
\def\labelenumi{\arabic{enumi}.}
\item
  \textbf{Blocking --- boundary and bulk symmetries confused (95--105,
  618--620).} Boundary global symmetries can correspond to bulk gauge
  symmetries. Flagging the Ising global symmetry as a gravity
  inconsistency reverses the relevant distinction.
\item
  \textbf{Blocking --- anomaly logic (101--115).} A nontrivial 't Hooft
  anomaly can obstruct gauging without making the original theory
  inconsistent.
  \ensuremath{\Omega}\ensuremath{{}_{2}}\textasciicircum{}Spin=Z\ensuremath{{}_{2}}
  does not itself say that every nontrivial Arf theory has the asserted
  gravitational anomaly. The rewrite confines the check to explicit
  symmetry data.
\item
  \textbf{Major --- rational versus algebraic arithmetic (88--93,
  613--615).} Modular matrices need algebraic numbers, not only
  rationals. Finite modular matrices also presuppose a suitable rational
  chiral algebra.
\item
  \textbf{Major --- unbounded classification and duplicates.} A
  central-charge ceiling alone does not define a finite RCFT search. The
  two 08 descriptions are consolidated, and higher-form/cobordism
  extensions require their spacetime dimension, symmetry type, and
  inflow theory before computation.
\end{enumerate}

\subsubsection{What the revision preserves and changes}

The central motivation is retained, but the executable question is now:
\textbf{For supplied rational-CFT modular data and a specified finite
internal symmetry, which algebraic consistency tests can be certified,
and which failures obstruct a proposed gauging?} The revised scope
chooses a concrete starting model; those choices are proposed
research-design decisions, not facts inferred from the original. The
baseline can succeed without resolving the full open research program.

The replacement also removes the original implementation sketches from
the active specification. They remain in the archive for comparison, but
have not been made executable or certified. New implementations should
derive their equations from the revised definitions rather than copy
unchecked prototypes.

\subsubsection{Acceptance criteria to use instead}

Baseline: exact modular-data checker on Ising and one second supplied
model. Strong: one finite-group twisted-sector/gauging analysis.
Research extension: classify a explicitly finite family of data or
derive a conditional bulk consequence using a stated holographic
theorem.

\subsubsection{Sources supporting the corrections}

\begin{itemize}
\tightlist
\item
  \href{https://arxiv.org/abs/1810.05338}{Harlow and Ooguri, Symmetries
  in quantum field theory and quantum gravity} --- Bulk gauge versus
  boundary global symmetry, including higher-form extensions.
\end{itemize}

\section{Evidence and execution standard}

This standard applies to the revised specifications. It is a proposed
contract for future research runs; this editorial revision does not
supply the scientific certificates themselves.

\subsection{What ``pure thought'' means here}

Use explicitly stated mathematical models and deductions without fitting
conclusions to experimental outcomes. Synthetic examples, exact model
inputs, analytic references, numerical search, and established software
are allowed. Physical constants, electronic-structure approximations,
constitutive laws, and empirical parameter sets must be identified when
used. A model-derived result is conditional on that model.

The original ban on legacy software conflicts with its proposed reliance
on NumPy, SciPy, CVXPY, nauty, and other libraries. Reimplementing every
optimizer and algebra system is not necessary for scientific
independence. Prefer a small independently checkable certificate and a
documented trusted base. A fresh implementation can be a project
objective, but must not be confused with a mathematical proof or novelty
claim.

\subsection{Four evidence levels}

{\def\LTcaptype{none} % do not increment counter
\begin{longtable}[]{@{}
  >{\raggedright\arraybackslash}p{(\linewidth - 4\tabcolsep) * \real{0.3333}}
  >{\raggedright\arraybackslash}p{(\linewidth - 4\tabcolsep) * \real{0.3333}}
  >{\raggedright\arraybackslash}p{(\linewidth - 4\tabcolsep) * \real{0.3333}}@{}}
\toprule\noalign{}
\begin{minipage}[b]{\linewidth}\raggedright
Label
\end{minipage} & \begin{minipage}[b]{\linewidth}\raggedright
What it establishes
\end{minipage} & \begin{minipage}[b]{\linewidth}\raggedright
What it does not establish
\end{minipage} \\
\midrule\noalign{}
\endhead
\bottomrule\noalign{}
\endlastfoot
Numerical evidence & A reproducible computation at stated precision,
sample size, and discretization & Exact truth, continuum coverage, or a
universal theorem \\
Finite-model certificate & An exact/validated statement about the
specified finite model or finite relaxation & An omitted tail,
thermodynamic limit, or physical realization \\
Conditional mathematical theorem & A proof over the stated domain with
every hypothesis and error controlled & Truth outside those assumptions
or a demonstrated physical system \\
Formal proof & A checked theorem in a named proof assistant, with
imported axioms/dependencies listed & Correct physical modeling merely
by virtue of formalization \\
\end{longtable}
}

Formalization is an optional extra deliverable unless the specific
project makes it mandatory. High precision, JSON output, solver status
and small residuals do not by themselves confer any of these proof
levels.

\subsection{Claim directions}

For minimization over a larger feasible set, the optimum is a lower
bound on the original minimum. A physical feasible point supplies an
upper bound. For maximization the directions reverse. Track feasibility
of the primal/dual witnesses separately from their objective values.

Infeasibility of an inner ansatz excludes only that ansatz.
Infeasibility of an outer relaxation can exclude the original problem if
the map into the relaxation is proved and all approximations preserve
necessity. Passing necessary tests means ``not excluded''; existence
needs a constructive sufficient argument.

A smallest observed example is not a minimality theorem. A complete
catalog must define the universe, equivalence relation, boundaries, and
proof that pruning misses nothing. A finite-size numerical trend does
not prove an asymptotic limit.

\subsection{Minimum run manifest}

Before calculation, record the problem ID and revision hash, exact
inputs and conventions, quantified domain, equivalence relation, target
observable, accepted evidence level, tolerances, and resource limits.
Record dependencies and precision/rounding mode. Where randomization is
used, preserve seeds, trial counts and raw sufficient statistics.
Resource limits are planning decisions, not evidence that a mathematical
solution cannot exist.

An output bundle should include:

\begin{itemize}
\tightlist
\item
  The statement being claimed, including all assumptions and
  quantifiers.
\item
  Construction or proof object, plus an independently runnable verifier
  and its trust assumptions.
\item
  Numerical/error ledger separating roundoff, discretization,
  truncation, optimization, sampling and modeling error where relevant.
\item
  Positive reference cases and negative controls that the verifier must
  reject.
\item
  Coverage: included domain, omitted domain, unresolved regions and
  termination reason.
\end{itemize}

These are requirements for future implementations, not a promise of
polynomial-size certificates for every problem.

\subsection{Checking specific certificate types}

\textbf{LP/SDP/SOS:} verify equality constraints, objective direction,
and cone feasibility using exact arithmetic or outward-rounded
intervals. Rationalizing a numerical solution can break PSD constraints.
Correct residuals with justified norm bounds; do not silently project
and reuse the old objective. Distinguish a proof of a polynomial
approximation from a proof for the original function.

\textbf{Interval numerics:} state function enclosures, domains and
outward rounding. A small root box needs existence and, if claimed,
uniqueness. Cover the complement for a global enumeration. Include
singularities and tails rather than discarding them numerically.

\textbf{SAT/SMT:} a solver's UNSAT label is not a certificate. Preserve
a proof supported by a checker and the proof that the encoding, symmetry
reduction and blocked solutions match the mathematical universe. DRAT
applies to an appropriate Boolean CNF proof, not automatically to every
real-arithmetic or SDP claim.

\textbf{Statistics:} specify the estimand, dependence assumptions,
confidence procedure and bias. Zero observed failures does not prove
zero failure probability. MCMC error bars based on estimated
autocorrelation are not generally rigorous mixing guarantees.

\subsection{Milestones and novelty}

Baseline means a working, independently validated result on a concrete
reference problem. Strong means a broader or sharper result with its
claimed evidence level. Research extension means a possible novel result
whose feasibility and significance need review. A negative theorem can
be valuable; an unsuccessful search can also be useful if its scope is
documented, but the two are not interchangeable.

Do not promise ``publication-quality'' in a fixed number of months.
Before claiming novelty, conduct a focused current literature comparison
against the exact theorem/model/metric. The sources in this revision
check key definitions and limitations; they are not an exhaustive
novelty audit across all 30 fields.

\end{document}
